\documentclass[12pt]{article}
\usepackage{amsmath}
\usepackage[margin=1in]{geometry}

\title{Labwork 3: Logistic Regression}
\author{Trung Tran Minh - 2540024}
\date{}

\begin{document}
\maketitle

\section{Introduction}

In this lab, I implement logistic regression from scratch by using gradient descent.
The input is a CSV file with three columns. The first two columns are input features and the last column is the label.
In this dataset $Loan.csv$ within features are Experience and Salary.

\section{Result}

I use learning rate:

\[
r = 0.1
\]

and run gradient descent for 10000 iterations.
The final result is:

\[
w_0 = -12.368200
\]

\[
w_1 = 4.249180
\]

\[
w_2 = 1.147938
\]

The final loss is:

\[
J = 0.210407
\]

The accuracy on this small dataset is:

\[
Accuracy = 0.8947
\]

So the logistic regression model is:

\[
\hat{y} = \frac{1}{1 + e^{-(w_0 + w_1 + w_2)}}
\]

\section{Effect of Different Learning Rates}

I also try different learning rates. The result is:

\[
\begin{array}{c|c|c|c|c|c}
r & w_0 & w_1 & w_2 & loss & accuracy \\
\hline
0.001 & -0.3310 & 1.2495 & -0.1400 & 0.4734 & 0.7368 \\
0.01 & -3.6249 & 2.2565 & 0.1946 & 0.3313 & 0.7895 \\
0.1 & -12.3682 & 4.2492 & 1.1479 & 0.2104 & 0.8947 \\
1.0 & -28.6238 & 8.5510 & 3.1459 & 0.3233 & 0.9474
\end{array}
\]

From this table, a very small learning rate learns slowly.
A learning rate of $0.1$ gives a good result for this dataset.
When the learning rate is too large, the result can be unstable and the loss does not always become smaller.

\end{document}
